% Part: incompleteness
% Chapter: theories-computability
% Section: first-incompleteness

\documentclass[../../../include/open-logic-section]{subfiles}

\begin{document}

\olfileid{inc}{tcp}{inc}

\olsection{$\Th{Q}$ دی کوئی مکمل، بے تضاد، !!^{axiomatizable}
  توسیع نہیں}

\begin{thm}
\ollabel{thm:first-incompleteness}
$\Th{Q}$ دی کوئی مکمل، بے تضاد، !!{axiomatizable} توسیع نہیں۔
\end{thm}

\begin{proof}
اسیں پہلاں ایہ ثابت کر چکے آں کہ $\Th{Q}$ دی کوئی بے تضاد
توسیع قابلِ فیصلہ نہیں۔ پر جے $\Th{T}$ مکمل تے !!{axiomatized}
اے، تاں اوہ قابلِ فیصلہ اے۔
\end{proof}

\begin{explain}
ایہ قضیہ گوڈل دے 1931 والے پہلے نامکملیت قضیے دی اصل صورت
دے نیڑے اے۔ نویں اصطلاحاں توں چھٹ، دو وڈے فرق نیں۔ گوڈل دی
شرط عام بے تضادی دی تھاں $\omega$-بے تضادی سی۔ اوہ “!!{axiomatizable}”
نوں پورے عموم نال وی نہیں ورت سکدا سی، کیوں جے قابلِ حسابیت دا
رسمی تصور اجے نہیں سی بنیا۔ (قابلِ حسابیت دے رسمی نمونے اگلے
دس ورھیاں وچ بنے۔ گوڈل نے وی ایس کم وچ حصہ پایا۔ ایہ نمونے
بنان دا اک وڈا مقصد ایہ پہچاننا سی کہ کہڑے نظرییاں اُتے گوڈل
دا نتیجہ لاگو ہو سکدا اے۔)

قضیہ آکھدا اے کہ توسیع وچ مکمل ہون، بے تضادی تے !!{axiomatizability}
تینوں کٹھیاں نہیں ہو سکدیاں۔ پر جے کوئی اک شرط چھڈ دیو، تاں
باقی دو پوریاں ہو سکدیاں نیں۔ $\Th{Q}$ بے تضاد اے تے اوہدی
مسلّمہ بندی قابلِ حساب اے، پر اوہ مکمل نہیں۔ تضاد والا نظریہ
مکمل اے تے اوہدی مسلّمہ بندی قابلِ حساب اے (مثال لئی
$\{ \eq/[0][0] \}$ نال)، پر اوہ بے تضاد نہیں۔ حساب دے سچے
جملیاں دا سیٹ مکمل تے بے تضاد اے، پر اوہدی مسلّمہ بندی قابلِ
حساب نہیں۔
\end{explain}
\end{document}
